\thispagestyle{empty}
\begin{tcolorbox}[enhanced, colback=hblue, colframe=hblue, boxrule=0pt, arc=4pt, left=16pt, right=16pt, top=22pt, bottom=22pt]
  \color{white}\sffamily
  {\large\bfseries Infinitesimaalrekening \textperiodcentered\ Blok 1 \textperiodcentered\ 2026--2027}\\[10pt]
  {\fontsize{34}{38}\selectfont\bfseries Workbook}\\[6pt]
  {\fontsize{20}{24}\selectfont Chapters 0--7}\\[14pt]
  {\normalsize Minimal theory, taken from the dictaat's own definitions, theorems and strategies, followed immediately by the dictaat's own exercises: every Oefening of Hoofdstuk 0--7 (1--30, 101--123, 201--220, 301--330, 401--420, 501--524, 601--635, 701--730) is either worked in the text or solved in Appendix S.}
\end{tcolorbox}

\vspace{10pt}
\begin{tcolorbox}[enhanced, colback=white, colframe=hblue2, boxrule=0.8pt, arc=2pt, left=10pt, right=10pt, top=8pt, bottom=8pt]
\sffamily\small\sloppy
\begin{itemize}[leftmargin=1.2em]
\item \textbf{\textcolor{hblue}{Source.}} S.~Wepster, \emph{Infinitesimaalrekening}, editie 2026, Mathematisch Instituut, Universiteit Utrecht: Blok 1, Hoofdstuk 0--7 (pages 17--113). Definition, Stelling, Voorbeeld and Oefening numbers are the dictaat's; nothing is used that the dictaat does not contain or take for granted (vwo algebra and, in the marked exercises, basic calculus).
\item \textbf{\textcolor{hblue}{Which exercises are worked where.}} Roughly half of the exercises, chosen to show each technique once, are worked in full in the chapter text. All others appear in NOW YOU boxes and are solved in Appendix S. Exercises the dictaat marks \emph{(Extra)}, and the (Extra) sections 3.5, 6.6 and 7.1, are outside the exam material but included.
\item \textbf{\textcolor{hblue}{How to use it.}} Read one theory box, cover the solution of the worked exercise that follows, attempt it, then compare. Do the NOW YOU exercises before opening Appendix S. Reflect after every exercise, as the dictaat asks: could you have seen the answer coming, would another approach work, what if the question were changed slightly?
\end{itemize}
\end{tcolorbox}

\vspace{8pt}
\hsub{How this workbook works}
Every section is a sequence of short cycles: \textbf{one piece of theory, then the exercises that use it.}

\begin{center}\sffamily\small
\begin{tabular}{@{}p{0.30\linewidth}p{0.65\linewidth}@{}}
\colorbox{hyellow}{\strut\textbf{Definitie n.m}} & The dictaat's definitions and notation, condensed; the things it lists under ``hoofdpunten: uit je hoofd weten''.\\[3pt]
\colorbox{hlight}{\strut\textbf{Stelling / Voorbeeld}} & A result you may use, with the idea behind it, or a worked example from the dictaat that serves as the model for the exercises.\\[3pt]
\colorbox{hgreen}{\strut\textbf{Strategy}} & The dictaat's methods (separating variables, integrating factor, substitution, partial fractions, function investigation) as step-by-step recipes.\\[3pt]
\colorbox{white}{\strut\textcolor{hblue}{\textbf{Exercise n}} \textcolor{horange}{\small\textbf{badge}}} & An Oefening worked in full, with the checks the dictaat insists on (differentiate your primitive, substitute your solution, sanity-check a value). Cover the solution and try first.\\[3pt]
\colorbox{hpurplebg}{\strut\textbf{NOW YOU}} & The remaining exercises for that piece of theory, all solved in Appendix S.\\[3pt]
\colorbox{hredbg}{\strut\textbf{Watch out}} & Where marks are lost: the dictaat's own warnings and the errors its author sees most in exams.
\end{tabular}
\end{center}

\begin{key}[{Conventions of the dictaat}]
\begin{itemize}
\item $\log$ is the \textbf{natural} logarithm (base $e$); $\ln$ is never used. $\N=\{0,1,2,\dots\}$.
\item The inverse function of $f$ is written $f_{\mathrm{inv}}$, not $f^{-1}$ (to avoid confusion with $1/f$); likewise $\arcsin,\arccos,\arctan$ and $\operatorname{arsinh},\operatorname{arcosh},\operatorname{artanh}$.
\item ``The function $f(x)$'' is sloppy but used; $f:A\to B$, $x\mapsto f(x)$ is correct. Intervals $(a,b)$ open, $[a,b]$ closed; $\infty$ is never a number.
\item Exercises are numbered per chapter: Oefening 203 is the third exercise of Hoofdstuk 2. The dictaat marks easy theory-checks and ``switching'' problems with icons that are not reproduced here; (Extra) marks exercises outside the exam material.
\item Exam: no calculator, no formula sheet; sum and double-angle formulas (Oefening 26) and the standard primitives (\S6.1) must be known by heart.
\end{itemize}
\end{key}
